\input{\string~/Documents/School/"UC Irvine"/ta_template2.0.tex}
\rh{2B}{7.1-7.8}
\chead{\today}
\graphicspath{ {\string~/Desktop} }
\usetikzlibrary{arrows}
%============ANSWER KEY TOGGLE===========
\newcommand{\ans}[1]{ \noindent {\color{red} \textbf{Answer:} #1}}
%\renewcommand{\ans}[1]{\vspace{3in}}
%==========================================
\begin{document}
\begin{center}
\textbf{MATH2B Midterm II Review}
\end{center}

\begin{ins}
The key to this part of the course is mastering a few procedures dealing with integrals, and recognizing when to use them. Below is an itemized list of the sections covered since the last exam, each with a collection of associated problems. This table of contents should also serve as a guide to thought process: when faced with a new integral, use this as a hierarchy of methods to try. 
\end{ins}

\begin{aun}
This document has \emph{a lot} of problems. If you have any questions, feel free to \href{mailto:vinnakot@uci.edu}{email me.}
\end{aun}


\tableofcontents


\section{$u$-Substitution}
\begin{aun}
Thing to remember: Look for a composition of functions, and differentiate the inner function.
\end{aun}

\problembox{When integrating $\int x^3\left(x^4+1\right)^2 d x$ using the substitution method, we would begin by letting $u$ equal:
\begin{enumerate}[label=\alph*.]
\item $x^3$
\item $x\left(x^4+1\right)$
\item $\left(x^4+1\right)^2$
\item $x^4+1$
\end{enumerate}
}
\ans{Inner function here is $x^4 + 1$. }

\problembox{Evaluate the integral:
\[
\int \frac{\lrp{\ln(x)}^6}{x}dx
\]}
\ans{ $\dst \frac17 \ln(x)^7 + C$ }

\problembox{Using the substitution $u=x^2-3, \int_{-1}^4 x\left(x^2-3\right)^5 d x$, what is equal to?}

\ans{$\dst \int_{-2}^{13} u^5 du$}


\section{Integration by Parts}
\begin{aun}
Thing to remember:
\[
\int u dv = uv - \int v du
\]
\end{aun}

\begin{aun}
\href{https://www.math.ucdavis.edu/~kouba/CalcTwoDIRECTORY/intbypartsdirectory/IntByParts.html}{Here's a link to review}
\end{aun}


\problembox{Compute
\[
\int x \cos(x) dx
\]
}
\ans{$x \sin (x) + \cos (x) + C$}

\problembox{Compute
\[
\int (\ln(x))^2 dx
\]
}
\ans{Do this by $u$-substitution first: let $u = \ln(x)$ and $du = 1/x dx$
\[
	\int u^2 (e^u du)
\]
Do integration by parts twice to get:
\[
x ((\ln(x))^2-2 \ln(x)+2) + C
\]
 }

\problembox{
Evaluate the definite integral:
\[
\int_0^1 x^2 e^x dx
\]
}
\ans{$e-2$}

\problembox{
Let $f$ be a differentiable function such that $\int f(x) \sin x d x=-f(x) \cos x+\int 4 x^3 \cos x d x$. What is a possibility for $f(x)$?
}

\ans{$x^4$}

\section{Trigonometric Integration}
\begin{aun}
Thing to remember:
\[
\sin^2 (x) + \cos^2(x) = 1 \mte{ and } \tan^2 (x) + 1 = \sec^2 (x)
\]
Look for \emph{odd} powers on $\sin, \cos$ or factors of $\sec^2$ and $\sec \tan$. 
\end{aun}
\begin{aun}
\href{https://www.math.ucdavis.edu/~kouba/CalcTwoDIRECTORY/trigintdirectory/TrigInt.html}{Here's a link to review}, \href{https://tutorial.math.lamar.edu/classes/calcii/integralswithtrig.aspx}{here's another}
\end{aun}

\problembox{Integrate:
\[\int (\sin(x) + \cos(x))^2 dx\]
}

\ans{
Expand out the polynomial to get:
\[
	\underbrace{\sin^2(x) + \cos^2(x)}_{=1} + 2 \sin(x) \cos(x)
\]
The rest should be just $u$-substitution:
\[
x - 2\sin^2(x) + C
\]
}


\problembox{Compute:
\[
\int\left(\sec ^2 x\right) \sqrt{5+\tan x} d x
\]
}
\ans{Eyes should light up at the $\sec^2(x)$ term hanging out there, this is just a $u$-sub on $u = \tan(x)$. 
\[
= \frac23 (5 + \tan(x))^{3/2} + C
\]
}

\section{Trigonometric Substitution}
\begin{aun}
Thing to remember:
\[
\begin{array}{|c|c|c|}
\hline
\text{Expression} & \text{Substitution} & \text{Identity} \\ \hline\hline 
\sqrt{a^2 - x^2} & \displaystyle x = a \sin \theta, \theta \in \lrb{\frac{-\pi}{2}, \frac\pi2} & 1 - \sin^2 \theta = \cos^2 \theta\\ \hline 
\sqrt{a^2 + x^2} & \displaystyle x = a \tan \theta, \theta \in \lrp{\frac{-\pi}{2}, \frac\pi2} & 1 + \tan^2 \theta = \sec^2 \theta \\ \hline
\sqrt{x^2 - a^2} & \displaystyle x = a \sec\theta,  \theta \in \left [ 0, \frac\pi2 \right ) \text{ or } \in \left [ \pi, \frac{3\pi}{2} \right ) & \sec^2\theta - 1 = \tan^2 \theta \\ 
\hline 
\end{array}
\]
To complete problems, use a triangle.
\end{aun}

\begin{aun}
\href{https://tutorial.math.lamar.edu/classes/calcii/TrigSubstitutions.aspx}{Here's a link to review} \href{https://www.math.ucdavis.edu/~kouba/CalcTwoDIRECTORY/trigsubdirectory/TrigSub.html}{and another}
\end{aun}

\problembox{Compute
\[
\int \sqrt{1 - x^2} dx
\]
}
\ans{Deceptively tricky! Make the substitution $x = \sin(\theta)$ to get the integrand:
\[
\int \cos^2(\theta) = \frac12 \int 1 + \cos(2\theta) d \theta
\]
with the latter equality from the \emph{power reduction} formula. Integrating gives:
\[
	\frac12 \theta + \frac14 \sin(2\theta) = \frac12 \theta + \frac12 \sin (\theta) \cos (\theta )
\]
Reversing the subsitution using the Pythagorean theorem gives:
\[
		\boxed{\frac12 \arcsin(x) + \frac12 x \sqrt{1-x^2} + C}
\]
}

\problembox{Compute
\[
\int_0^1 \sqrt{x^2-2 x+1} d x
\]}

\ans{$\frac12$}

\problembox{Compute
\[
\int_0^1\left(4-x^2\right)^{-\frac{3}{2}} d x=
\]}

\ans{Make the subsitution $x=2 \sin \theta \Rightarrow d x=2 \cos \theta d \theta$.
$$
\int_0^1\left(4-x^2\right)^{-3 / 2} d x=\int_0^{\pi / 6} \frac{2 \cos \theta}{8 \cos ^3 \theta} d \theta=\frac{1}{4} \int_0^{\pi / 6} \sec ^2 \theta d \theta=\left.\frac{1}{4} \tan \theta\right|_0 ^{\pi / 6}=\frac{1}{4} \cdot \frac{\sqrt{3}}{3}=\frac{\sqrt{3}}{12}
$$}

\section{Partial Fraction Decomposition}
\begin{aun}
Thing to remember: decompose:
\[
\frac{p(x)}{q(x)} = s(x) + \frac{r(x)}{q(x)}
\]
using polynomial long division. Then remember: ``constants for distinct linear factors and their repetitions, linear terms for the quadratics''
\end{aun}
\begin{aun}
\href{https://www.math.ucdavis.edu/~kouba/CalcTwoDIRECTORY/partialfracdirectory/PartialFrac.html}{Here's a link to review}
\end{aun}

\problembox{Integrate:
\[
\int \frac{2x + 3}{x^2 - 9}
\]}
\ans{
	\[
		\frac12 \lrp{ 3\ln (3-x) + \ln (x+3)}
	\]
}

\problembox{Compute:
\[
\int \frac{1}{x^2 - 6x + 8} dx
\]}
\ans{ 
\[
\frac12 \lrp{\ln \lrabs{x - 4} - \ln \lrabs{x - 2}}
\]}

\problembox{Compute the following: \[\int_0^1 \frac{5 x+8}{x^2+3 x+2} d x\]}
\ans{$\ln(18)$}

\section{Improper Integrals}
\begin{aun}
Thing to remember: does the integral have an infinite bound? if so, replace with $b$ and evaluate the limit $\lim_{b \to \infty}$. Does the function have discontinuities on the domain of integration? If so, approach with limits from either side. If any limit diverges, the integral diverges.
\end{aun}

\problembox{
Evaluate the integral:
\[
\int_1^\infty e^{-3x}dx
\]
}
\ans{$\dst \frac{1}{3e^3}$}

\problembox{Compute the integral
\[
\int_0^\infty x^2 e^{-x^3} dx
\]
}

\ans{
\[
\lim_{b \to \infty} \lrb{-\frac13 e^{-x^3} }^b_0 = \frac13
\]
}


\section{The Grab Bag}
\begin{aun}
If I sort these into a category, it'll give away the answer
\end{aun}

\problembox{[True or False] Integrals of the form:
\[
\int_0^\infty f(x)dx
\]
have a finite numerical value.
}
\ans{False}


\problembox{Evaluate the integral:
\[
\int\left(4 x^2-1\right) \cos \left(4 x^3-3 x\right) d x
\]
}
\ans{$=-\frac{1}{3} \sin \left(3 x-4 x^3\right)+\text { constant }$}


\problembox{Compute the integral:
\[
\int_0^1 \frac{x^2}{x^2 + 1}
\]
\textbf{HINT:} what's the derivative of $\arctan$?
}
\ans{$1 - \pi/4$}

\problembox{Compute \[\int_1^e\left(\frac{x^2-1}{x}\right) d x=\]}
\ans{$\frac12 (e^2 - 3)$; split this one into two integrals}

\problembox{If $\frac{dy}{dx} = \sin(x) \cos^2 (x)dx$ and $y = 0$ when $x = \frac{\pi}{2}$, what is $y$ when $x$ is 0?}

\ans{
Using the antiderivative 
\[
y(x) = -\frac13 \cos^3 (x) + 0 \mte{ we have }-\frac13
\]
}

\problembox{Compute
\[
\int_0^1 \sqrt{x} (x+1) dx
\]
}
\ans{ $\frac{16}{15}$}

\problembox{
Compute
\[\int_1^{\infty} \frac{x}{\left(1+x^2\right)^2} d x\]
}
\ans{
\[
\lim_{b \to \infty} \lrb{\frac14 - \frac{1}{2(1-b^2)} } = \frac14
\]
}

\problembox{Compute 
\[
\int \sin^3 (x)dx
\]
}
\ans{ Odd power of $\sin$ is good news, at least!
\[
\int( 1-\cos^2 (x)) \sin(x)dx = \int u^2 - 1 du ...
\]
Follow out the algebra to get 
\[
		-\cos(x) + \frac13 \cos^3(x) + C
\]
}

\problembox{Compute
\[
\int \frac{\cos^2(x)}{1 + \sin(x)}
\]
}
\ans{Conjugate the denominator:
\[
	\int \frac{\cos^2(x) (1-\sin(x))}{1- \sin^2(x)} = \int \frac{\cos^2(x) (1-\sin(x))}{\cos^2(x)} = \int 1- \sin(x) dx
\]
Finish it out from here to get:
\[
	x + \cos(x) + C
\]
}

\problembox{Compute 
\[
\int \frac{d x}{(x-1)(x+3)}
\]}

\ans{ Use partial fractions
\[
\frac14 \ln \lrabs{\frac{x-1}{x+3}} + C
\]
}

\problembox{Compute \[\int_4^{\infty} \frac{-2 x}{\sqrt[3]{9-x^2}} d x\]}
\ans{DNE by comparison to $x^{2/3}$}

\problembox{Compute \[\int \frac{\left(x^2-1\right)^{3 / 2}}{x} d x\]}
\ans{The radical should suggest a trig substitution:
\[
\int \frac{\left(x^2-1\right)^{3 / 2}}{x} d x = \int \frac{(\sqrt{x^2 - 1})^3}{x} dx = \int \frac{\tan^3 \theta}{\sec \theta}
\]
over the substitution $x = \sec(\theta)$. 
\[
	\cdots = \frac13 \sqrt{x^2 - 1} (x^2 -4 ) + \arctan \lrp{\sqrt{x^2 - 1}} + C
\]
}


\problembox{Integrate \[\int \frac{x}{\sqrt{x^2+4 x+5}} d x\]}
\ans{Do a substitution after completing the square:
\[
	x^2 + 4x + 5 = (x+2)^2 + 1
\]
Get:
\[
	\sqrt{x^2+4 x+5} - 2\ln \lrabs{\sqrt{x^2+4 x+5} + x + 2} + C
\]
}


\problembox{Compute 
\[
\int x^3 \ln (5x) dx
\]
}

\ans{This is parts- get:
\[
\frac14 x^4 \ln \lrp{5x} - \frac{x^4}{16} + C
\]}

\problembox{Compute 
\[
\int x \sin (x) \cos (x)dx
\]
}
\ans{Tricky tricky- see that:
\[
\int x \sin(x) \cos(x) dx = \frac12 \int x \sin(2x)dx
\]
Apply parts from here to get:
\[
	\frac12 \lrp{- \frac12 x \cos(2x) + \frac14 \sin(2x)} + C
\]
}

\problembox{Integrate
\[
\int \frac{x^4 + x^3 + x^2 + 1}{x^2 + x -2}
\]
}
\ans{Start with polynomial long division:
\[
	\frac{x^4 + x^3 + x^2 + 1}{x^2 + x -2} = x^2 + 3 + \frac{-3x+7}{x^2 + x -2}
\]
Continue by partial fractions:
\[
	= \frac13 x^3 + 3x + \frac43 \ln \lrabs{x-1} - \frac{13}{3} \ln \lrabs {x+2} + C
\]
}


\end{document}
